\section{Introduction}
\label{sec:intro}
Multi-task benchmarks and leaderboards (MMLU, HELM, the Open LLM Leaderboard) rank models by aggregating scores on many sub-datasets, typically with a mean. A leaderboard position influences adoption, publication and funding, so developers have a reason to act strategically, for instance by putting data related to some sub-datasets into a specific model's training in order to raise its scores there. Prior work shows that training on benchmark-related data can inflate scores \citep{zhou2023cheater,yang2023rephrased}, that leaderboard practices can distort rankings \citep{singh2025illusion}, and that incentives of this kind can be analysed with game-theoretic models \citep{chen2026incentives}. What is missing is a direct, computable answer to a concrete question: \emph{on a given public leaderboard, how much targeted investment does it take for a chosen model to become first, and how does that depend on the aggregation rule?}

We answer it with an explicit model (Section~\ref{sec:setting}): the target may raise its score on any task up to a per-task cap while all rivals stay fixed, and ``first'' means strictly above every rival by a small margin. We measure investment by the number of tasks touched $K$, the total points $L_1$ and the largest single-task rise $L_\infty$, each minimized exactly, and treat the cap as a scenario, never as an estimate of achievable training gains. We compare five score rules (mean, median, minimum, geometric and harmonic means) and three ordinal rules (Borda, plurality, Condorcet winner) on every model of two public leaderboards (score rules on all models; ordinal rules on HELM and a fixed top-500 frame of the Open LLM Leaderboard, with the Condorcet requirement an upper bound for Condorcet-consistent rules).

\paragraph{Contributions and scope of the claims.}
\begin{enumerate}\setlength\itemsep{0pt}
\item \emph{Theory (elementary, with proofs in Appendix~\ref{app:proofs}).} Closed forms for the mean, minimum and median (the cardinal analogues of margin-of-victory computations), a threshold equivalence between $L_\infty$ and per-task gain caps (for any rule with a closed first-placing set), and constructed examples proving that \textbf{no universal ordering of the mean, median and minimum rules holds} by $K$ or $L_1$ (Section~\ref{sec:setting}). The empirical census below is the main contribution.
\item \emph{Measurement.} Exact, solver-validated costs to be first, with ranges from resampling the targets, for all models of HELM MMLU (79 models, 57 subjects) and the Open LLM Leaderboard (4{,}496 models, 6 tasks), with infeasibility reported explicitly (Section~\ref{sec:experiments}). The answer is conditional on frame, cap and measure: raw points needed are largest under the mean; the number of tasks needed and, under a cap, whether first place is reachable at all, depend on the frame, the cap and the unit.
\item \emph{Controls and sensitivity.} Targeting beats random task order and uniform improvement, but a lowest-score-first greedy order already attains the optimal task count at the ceiling cap, so the optimiser's value lies in points and in capped settings; matched rival and task counts, task resampling and score noise show which comparisons are stable (Section~\ref{sec:experiments}, Appendix).
\end{enumerate}
All numbers are cost-to-first computations on fixed score snapshots; they describe those censuses and say nothing about what training on related data would actually deliver, nor about any developer's behaviour.
